\section{Preparing a Congestion Response}
The method combines an executable process model with learned behavior, compares feasible interventions under specified stresses, and evaluates the saved response on independent scenarios. Its output is an action with an observable trigger and a record of the assumptions under which it was selected.

\subsection{Represent execution and behavioral feedback}
A case passes through operations described by $\mathcal M=(N,\mathcal R,K,\alpha,F_D,F_R)$: control flow $N$, resource pools $\mathcal R$, capacities $K$, activity-to-pool mapping $\alpha$, and duration and routing functions. The supported subclass comprises block-structured sequences, XOR choices, AND splits and joins, and nested loops with visible bodies. Operations reserve input tokens when queued and produce output tokens on completion. Ready operations share finite FIFO pools, with occupancy $0\leq b_r(t)\leq K_r$. Work in progress $w(t)$ counts unfinished cases once, including cases with several simultaneous operations.

The lifecycle log records admission, enabled, start and completion times, activity, resource pool and observed case attributes. Duration is sampled at service start $t_s$:
\begin{equation}\label{eq:duration}
\log D_a=f_a(x,\ell(t_s))+\epsilon_a,\qquad \epsilon_a\sim\mathcal N(0,\sigma_a^2).
\end{equation}
At a binary return decision made at completion $t_c$,
\begin{equation}\label{eq:return}
R\sim\mathrm{Bernoulli}\!\left(\operatorname{sigmoid}(g(x,\ell(t_c)))\right).
\end{equation}
Here $x$ contains observed class and, for competing processes, process identity. Load features include $\log(1+q_a/K_{\alpha(a)})$, $\log(1+w/S)$ and process-specific WIP. The scale $S$ is fixed before evaluation. The start queue includes the operation about to start; completion-time features exclude a prospective return. Future states are never input features.

For multiple routes, a per-source multinomial model predicts the next enabled activity or completion of the case. Simultaneously enabled activities use a fixed lexical convention in the training target. Each executable choice receives the summed predicted probability of its reachable next labels; binary loop probabilities are then normalized. Class context distinguishes heterogeneous routes, while separate source models distinguish different return points. Duration uses ridge regression with penalty $\alpha=1$; routing uses logistic regression with $C=10$. Load terms in routing apply at the operationally specified repeat-intervention sites.

Static models retain attributes and resource limits but omit load features. Retained-load models include both service and routing dependencies. Predictive selection keeps each component only when its held-out probabilistic-loss improvement exceeds $0.01$ for duration or $0.005$ for routing. The losses are Gaussian negative log likelihood for log duration and cross entropy for route outcomes. This comparison isolates whether predictive selection preserves information needed for a later response decision. The original serial implementation P0 retains its original feature scaling, coefficients and results; the structural extension uses separate activity and return-source models.

The structure can be supplied or recovered from training traces. Discovery compares standard Inductive Miner with a candidate that tries its strict-loop decomposition before permissive concurrency. The latter also infers binary guards on class or process identity, observed at admission and constant within a case, when each attribute value has at least 10 training cases with a constant, distinct activity set. Among executable candidates fitting every training trace under optimal alignments, selection maximizes training precision, then minimizes transition count. Resource mappings, candidate return sites, their queue-observation mappings and action semantics are engineering inputs; true control flow and test outcomes do not enter selection.

Execution must preserve the meaning of the intervention site. A next-activity model cannot distinguish alternatives that begin with the same visible continuation. For example, $\mathrm{seq}(\mathrm{loop}(B,D),D)$ and $\mathrm{loop}(\mathrm{seq}(B,D),\tau)$ admit the same visible traces, but the latter places the repeat decision after D, where its outcome and intervention effect can be represented. The compiler applies this normalization and rejects unresolved choices with overlapping next-label sets. Observed attribute guards are enforced during fitted simulation. The supported subclass therefore requires deterministic observed guards or distinguishable next-label alternatives after normalization. No discovered structure is repaired using the known generator.

\subsection{Compare feasible responses under specified stresses}
An engineer supplies a planning library $\mathcal S_p$, actions $\mathcal A$, costs and candidate WIP thresholds. A response $\pi=(a,q)$ activates after two successive threshold crossings:
\begin{equation}\label{eq:trigger}
\tau_q=\inf\{t_k:w(t_{k-1})\geq q\ \land\ w(t_k)\geq q\}.
\end{equation}
It uses observed WIP without knowledge of future arrivals. Adding a worker affects future starts, while reducing rework affects future return decisions. Already sampled service times remain unchanged. With horizon $H$, the objective is
\begin{equation}\label{eq:cost}
J(\pi)=\int_0^H w_\pi(t)\,dt+c_a(H-\tau_q)_+,
\end{equation}
with zero action cost if activation never occurs. The WIP integral uses trapezoidal integration of observations. Maintaining one additional worker costs $2$ per time unit; maintaining rework reduction costs $0.5$. These are case-time equivalents, not monetary estimates. Terminal unfinished work is reported separately.

The planner chooses the lowest mean simulated cost over specified scenarios and repetitions:
\begin{equation}\label{eq:policy}
\widehat\pi=\arg\min_{\pi\in\Pi}\frac{1}{|\mathcal S_p|L}
\sum_{s\in\mathcal S_p}\sum_{l=1}^{L}J_{\widehat{\mathcal M},s,l}(\pi).
\end{equation}
All model variants use the same actions, thresholds, scenarios and costs. Ties follow fixed catalog order. Counter-based random streams pair action comparisons. The selected policy is saved before independent test arrivals are generated.

\subsection{Evaluate the saved response}
The test measures realized cost, sustained-overload incidence, activation, lead time and completed-case duration. Lead time refers to the onset of severe overload in the matched no-action trajectory; that counterfactual is used only for evaluation. Preventing an event differs from reducing its accumulated cost. A response can lower the backlog while the severe-event criterion remains satisfied.

The response record connects the action and trigger to the fitted model, planning stresses, observed load range, intervention assumptions and test outcomes. These fields support operational approval and identify when new measurements or recalibration are required.
